% Source-derived chapter 09: Proofs of Theorems \ref{Adagio} and \ref{brahms}
% Original source: rigid-local-systems-multiplicative-eigenvalue-problem
\chapter{Proofs of Theorems \ref{Adagio} and \ref{brahms}}
\label{rigid-local-systems-multiplicative-eigenvalue-problem-chapter-09}
\source{rigid-local-systems-multiplicative-eigenvalue-problem}

\begin{remark}[Source section]
\label{rigid-local-systems-multiplicative-eigenvalue-problem-chapter-09-source}
\source{rigid-local-systems-multiplicative-eigenvalue-problem}
\source{rigid-local-systems-multiplicative-eigenvalue-problem}
This chapter reproduces the corresponding section of the retrieved
LaTeX source, including its definitions, statements, examples, and
proofs. It has not yet been translated into Lean.
\notready
\end{remark}

% The source section title was: Proofs of Theorems \ref{Adagio} and \ref{brahms}
\subsection{Shift operations}
The spaces ${\Omega}^0(d,r,\mn,n,\vec{I})$ (but not $\widehat{\Omega}^0(d,r,\mn,n,\vec{I})$) together with their maps to $\ParL$  behave well under shift operations:

Let $\deg\mn= -D$ and $p=p_i\in S$. Parallel to Diagram \eqref{basicD} we have a diagram
 \begin{equation}\label{basicDnew}
\source{rigid-local-systems-multiplicative-eigenvalue-problem}
\xymatrix{
 \Omega^0(d,r,\mn,n,\vec{I})\ar[d]\ar[r]^{\operatorname{Sh}_p}  & \Omega^0(d',r,\mn(p),n,\vec{K})\ar[d]\\
\ParL(-d,r,n-r,d-D)\ar[r]^{\operatorname{Sh}_p}  & \ParL(-d',r,n-r,d'-D+1)
}
\end{equation}
where the terms $d'$ and $\vec{K}$ are determined as in Diagram \eqref{basicD}:
 \begin{enumerate}
 \item $d'=d$ if $1\notin I^i$, and $d'=d-1$ if $1\in I^i$
 \item $K^k=I^k$ for $k\neq i$.
 \item $I^i =\{i^i_1-1,i^i_2-1,\dots, i^i_r-1\}$ if $1\notin I^i$, and $K^i =\{i^i_2-1,\dots, i^i_r-1,n\}$ if $1\in I^i$.
 \end{enumerate}
 The shift operation in the bottom row of \eqref{basicDnew} is shift on one of the factors and identity on the other. It is a shift on the rank $r$ vector bundle factor if $d'=d-1$ and on the  rank $n-r$ vector bundle if $d'=d$.
 \subsection{A scaling property}\label{FSP}
\source{rigid-local-systems-multiplicative-eigenvalue-problem}
We have a self-contained proof of a generalization of a conjecture of Fulton (known before, see \cite[Remark 8.5]{BKq}). In Remark \ref{Bee132}, we explain how the version in Proposition \ref{qTheorems} (2) follows from the following:
\begin{proposition}
\source{rigid-local-systems-multiplicative-eigenvalue-problem}
\notready\label{fulC} For all integers $n\geq 0$,
\source{rigid-local-systems-multiplicative-eigenvalue-problem}
$\dim H^0(\ParL,\mathcal{O}(n\mathcal{R}_L))=1$.
\end{proposition}
\begin{proof}
 Note that $\mathcal{R}_L$ is the pullback of $\mathcal{R}$ under $\ParL\to {\Omega}^0(d,r,\mn,n,\vec{I})$. We can therefore reduce to the case  $1\not\in I^i$ for all $i$ using diagrams \eqref{basicDnew} and \eqref{basicD}.
Under this assumption, a function on $\ParL-\mathcal{R}_L$ pulls back to a function on $\widehat{\Omega}^0-\mathcal{R}$ which embeds in $\Par_{n,\mn,S}$ with complement of codimension $\geq 2$ by Corollary \ref{carole}. Since $h^0(\Par_{n,\mn,S},\mathcal{O})=1$, the proof is complete.
\end{proof}
\begin{proposition}
\source{rigid-local-systems-multiplicative-eigenvalue-problem}
\notready\label{rigid}
\source{rigid-local-systems-multiplicative-eigenvalue-problem}
$H^0(\Omega^0-\mathcal{R},\mathcal{O})=\Bbb{C}$.
\end{proposition}
\begin{proof}
Let $f\in H^0(\Omega^0-\mathcal{R},\mathcal{O})$. By Proposition \ref{fulC}, the restriction of the function $f$ to $\ParL$ is a constant $c$. Given a point $p\in \Omega^0$, there is a mapping $g:\Bbb{A}^1\to \Omega^0$ which
\begin{itemize}
\item[-] Sends $1$ to $p$ and $0$ lies in the image of $\ParL$ and equals the image of $p$ under the map $\Omega^0\to \ParL$.
\item[-] The isomorphism class of $g(t)$ remains constant for $t\neq 0$.
\end{itemize}
The existence of such a mapping can be explained more easily in the general group situation. Let $P$ be the standard parabolic for the Grassmannian $\Gr(r,n)$. Let $L$ be the corresponding Levi subgroup. A point of $\Omega^0$ corresponds to a principal $P$-bundle with some extra structure at the fibers over points of $S$ (see e.g., Lemmas 3.3 and 3.4 in \cite{BKiers}). We now use the Levification operation in  \cite[Section 3.8]{BKq} to produce the family $g$.
\end{proof}
\begin{remark}
\source{rigid-local-systems-multiplicative-eigenvalue-problem}
\notready\label{Bee132}
\source{rigid-local-systems-multiplicative-eigenvalue-problem}
Assume $d=0$ in Proposition \ref{fulC} (but $D$ is arbitrary). By a basic computation (see Proposition \ref{basicC} below), $\dim H^0(\ParL,\mathcal{O}(\mathcal{R}_L))$ is a tensor product
$H^0({\Par}_{r,\mathcal{O}(-d),S},\ma)\otimes H^0({\Par}_{n-r,\mathcal{O}(D-d),S},\mb)$ where the line bundles $\ma$ and $\mb$ are in strange duality (Definition \ref{Shosta}), and hence the dimensions of their spaces of sections are the same. The formula for $\ma$ in Proposition \ref{basicC} shows that it includes all line bundles of grade zero on ${\Par}_{r,\mathcal{O},S}$. This implies Proposition \ref{qTheorems} (2).
\end{remark}

\subsection{Proof of Theorem  \ref{Adagio} and Theorem \ref{brahms}}\label{rumble}
\source{rigid-local-systems-multiplicative-eigenvalue-problem}
Most of the steps are similar to the proof of \cite[Theorem 1.9]{BHermit}. The proof will use shift operations and we will prove  a generalization to ${\Par}_{n,\mn,S}$. We start with a $\langle\sigma_{I^1}, \sigma_{I^2},\dots,\sigma_{I^s}\rangle_{d,D}=1$, where $D=-\deg(\mn)$.  The inequality \eqref{wallnew} is replaced by \eqref{wallneww}, so that $\mf$ is now given inside  $\Pic_{\Bbb{Q}}({\Par}_{n,\mn,S})$ by equality in \eqref{wallneww}.

We first deal with the case that $a\neq 1$. We will show that $\Omega^0(d,r,\mathcal{N},n,\vec{J})\subset \widehat{\Omega}^0(d,r,\mathcal{N},n,\vec{I})$ is not contained in $\mathcal{R}$, the ramification locus of $\widehat{\Omega}^0(d,r,\mathcal{N},n,\vec{I})\to  {\Par}_{n,\mn,S}$.

Choose a point $x=(\mv,\mw,\me,\gamma)\in\Omega^0(d,r,\mathcal{N},n,\vec{I})$ such that $\widehat{\Omega}^0(d,r,\mathcal{N},n,\vec{I})\to  {\Par}_{n,\mn,S}$ is etale at $x$, and then modify an appropriate element in the flag $E^{p_j}_{\bull}$, so that the point lies now in  $\Omega^0(d,r,\mathcal{N},n,\vec{J})$. The modification is the following: replace  $E^{p_j}_{a-1}$ by $E^{p_j}_{a-2}+\mv_{p_j}\cap E^{p_j}_{a}$. This way we get a new $\me'\in \Fl_S(\mw)$, and $y=(\mv,\mw,\me',\gamma)\in\Omega^0(d,r,\mathcal{N},n,\vec{J})$. The tangent space ``situation'' at $y\in \widehat{\Omega}^0(d,r,\mathcal{N},n,\vec{I})$ is the same as at $x$, since the element $E^{p_j}_{a-1}$ does not appear in the rank conditions defining $\widehat{\Omega}^0(d,r,\mathcal{N},n,\vec{I})$. Therefore $y\not\in\mathcal{R}$.

The birational morphism $\widehat{\Omega}^0(d,r,\mathcal{N},n,\vec{I})\to  {\Par}_{n,\mn,S}$ is therefore  etale at $y\in \widehat{\Omega}^0(d,r,\mathcal{N},n,\vec{I})$. Since $D(a,j)$ is the  image of the irreducible  $\Omega(d,r,\mathcal{N},n,\vec{J})$, it is irreducible and it coincides with the cycle theoretic image $\widehat{\Omega}^0(d,r,\mathcal{N},n,\vec{I})\to  {\Par}_{n,\mn,S}$. This proves part (1) of Theorem \ref{Adagio}. Proposition \ref{enumerative} now implies Theorem \ref{brahms}.

 Note that $D(a,j)$ lies in the complement of the open subset $\Omega^0(d,r,\mathcal{N},n,\vec{I})-\mathcal{R}$ of ${\Par}_{n,\mn,S}$, since otherwise the image $y\in \Omega^0(d,r,\mathcal{N},n,\vec{J})$ in ${\Par}_{n,\mn,S}$ will have a disconnected inverse image in $\widehat{\Omega}^0(d,r,\mathcal{N},n,\vec{I})$. Part (2) of Theorem \ref{Adagio} now follows from Proposition \ref{rigid}. Part (3) of Theorem \ref{Adagio} follows from parts (1) and (2) using a suitable generalization of \cite[Lemma 2.1]{BHermit}. In fact, part(3) is already covered by Lemma \ref{elgar}.

Finally for part (4) we want that the line bundle $\mathcal{O}(D(a,j))$ on ${\Par}_{n,\mn,S}$ lies on the face $\mathcal{F}$. The numerical equality in \eqref{wallneww} is then just the statement that the center of $L$ acts trivially on $\mathcal{O}(D(a,j))$ restricted to $\Par_L$, which follows from the fact that the restriction is effective: Using arguments of Ressayre \cite{R1} this is equivalent to the line bundle restricted to $\ParL$ being effective, which is certainly the case because $D(a,j)$ does not contain the image of $\ParL$. Note  $\ParL-\mathcal{R}_L$ does not meet $D(a,j)$ since $\ParL-\mathcal{R}_L\subseteq \Omega^0(d,r,\mathcal{N},n,\vec{I})-\mathcal{R}$, and $D(a,j)$ does not intersect the last set.  This proves (4).

Finally we can remove the assumption $a\neq 1$. Assume $a=1\in I^j$. Consider the shifted data
$(d',r,\mathcal{N}(p_j),n,\vec{K})$ associated to $(d,r,\mn,n,\vec{I})$, with $d'=d-1$, with the shift applied at $p_j$, so that $K^j=\{i^j_2-1,\dots, j^j_n -1,n\}$. We refer to Theorem \ref{Adagio} (also see \eqref{basicD}) with the $\vec{I}$ replaced by $\vec{K}$. We now consider $\tilde{D}(n,j)$ for the tuple $(d',r,\mathcal{N}(p_j),n,\vec{K})$ and the above reasoning applies to it.  Using the inverse of the shift operation at $p_j$, and properties of the shift operation, we obtain the desired statement about $D(1,j)$.\subsection{Proof of Corollary \ref{existence}}
$\mathcal{B}(\vec{\lambda},\ell)$ is a F-line bundle by Theorem \ref{brahms}. The corollary now follows from Theorem \ref{christmas} and Corollary \ref{katzkatz}.
\iffalse

 The numerical equality is the rigidity condition \eqref{eqRi}: The eigenvalues of the monodromy of the rank $\ell$ local system at the point $p_i$ are the numbers $\exp(2\pi\sqrt{-1}b/n)$ with multiplicity $c^b_i$ for $1\leq b\leq n-1$ and $1$ with multiplicity $\ell-\sum_{i=1}^{n-1} c^b_i$ (the transpose of an $n\times \ell$ matrix $\sum_{b=1}^{n-1} c_b\omega_b$ has rows of length $b$ repeated $c_b$ times, and length $0$ repeated $\ell-\sum c_b$ times).

\fi
